\documentclass[11pt,oneside]{article}
\usepackage[margin=1in]{geometry}
\usepackage{booktabs}
\usepackage{amsmath,amssymb,amsthm}
\usepackage[hidelinks]{hyperref}
\usepackage{url}
\usepackage{cite}

% Shorthand
\newcommand{\bxthree}{BX3}
\newcommand{\purpose}{\textit{Purpose}}
\newcommand{\bounds}{\textit{Bounds}}
\newcommand{\fact}{\textit{Fact}}

\begin{document}

\section{Background}\label{sec:rs-background}

The problem of accountable recursive agent spawning sits at the intersection of four research traditions: agent lifecycle management in multi-agent systems, accountability and provenance in distributed AI architectures, fixed versus mutable delegation chain architectures, and hierarchical task decomposition in AI systems. Each tradition has produced significant contributions, but each also contains a specific structural gap that the Recursive Spawning Protocol (RSP) is designed to fill. This section surveys all four areas, establishes exactly where those gaps lie, and situates the BX3 Framework as the minimal sufficient substrate for immutable parent chains.

% ---------------------------------------------------------------
\subsection{Agent Lifecycle Management in Multi-Agent Systems}\label{sec:lifecycle}

Modern AI agent frameworks have moved away from treating agent creation as a singular one-time initialization event. Instead, agent creation, assignment, delegation, revocation, and termination are understood as managed transitions within a structured lifecycle \cite{wooldridge1995,stone2000,ferber1996}. Stone and Veloso established that agent management is not a deployment detail but a core architectural concern determining both scalability and accountability properties \cite{stone2000}. Ferber and M\"uller formalized that agent organizations are dynamic structures whose topology evolves with task demands, establishing the conceptual foundation for recursive delegation \cite{ferber1996}. Padgham and Thiel extended this to group membership dynamics, showing that agent collectives require explicit lifecycle support for formation, maintenance, and dissolution \cite{padgham2001}.

AGENTHIVE introduces a composable multi-agent framework in which delegation is treated as a first-class primitive: any agent can spawn and coordinate sub-agents, enabling decentralized control without a central orchestrator \cite{agenthive2025}. Through recursive delegation, task-driven structures emerge dynamically---trees, forests, and star topologies---that adapt to open-ended goals without requiring a predefined agent graph. Wu et al. demonstrate that this dynamic, delegation-centric architecture yields measurable performance gains across diverse real-world domains \cite{agenthive2025}. The key architectural insight is that delegation is not merely one-hop task assignment but can be \emph{recursive}: an agent may be assigned a task to identify and assign sub-tasks to others, effectively delegating the act of delegation itself \cite{arxiv2602.11865}. Furthermore, this decomposition may occur not only before task assignment but also dynamically during execution, meaning the depth and shape of the delegation tree are not fully specified at root-task initialization \cite{arxiv2602.11865}.

AgentSpawn demonstrates that dynamic spawning in response to runtime complexity metrics enables superior completion rates on long-horizon tasks compared to static baselines, with approximately 34\% improvement on SWE-bench and 42\% reduction in memory overhead via selective task slicing \cite{agentspawn2026}. The framework transfers relevant context to offspring agents automatically through memory management protocols, addressing memory continuity across agent generations through runtime-driven spawning policies \cite{agentspawn2026}. Coherence protocols ensure coordinated modifications across multiple agents working on overlapping regions, reducing conflicts and preserving task integrity during concurrent execution \cite{agentspawn2026}. AgentSpawn's spawning decision algorithm formalizes this as a metacognitive awareness: ``Am I the right agent for this subtask?''---a question that complements training-time evolution with architectural adaptability at execution time \cite{agentspawn2026}.

The critical gap in current lifecycle management approaches is the treatment of the parent reference. In most production multi-agent frameworks---including CrewAI's hierarchical process mode \cite{crewai2024}, LangGraph's stateful graph runtimes \cite{langchain2023}, and AutoGen's multi-agent conversation framework \cite{autogen2023}---the parent-child relationship is maintained as mutable runtime state. In CrewAI's hierarchical process, the manager agent holds references to worker agents that are updated dynamically as orchestration progresses. In LangGraph, node edges are part of the graph state that can be rewritten during execution, meaning the delegation chain is a living data structure rather than a fixed record \cite{zhuge2024}. Zhuge et al.'s holistic evaluation of LLM-based multi-agent systems confirms that provenance reconstruction in stateful graph runtimes is a persistent pain point, particularly in long chains where context divergence compounds \cite{zhuge2024}.

The Human Delegation Provenance (HDP) protocol comes closest to addressing this gap, proposing a cryptographic token that records human authorization decisions and propagates them through the delegation chain as first-class artifacts \cite{hdp2025ietf,hdp2025arxiv}. HDP's core insight is that human authorization is not a single permission granted at chain initialization but a per-decision artifact that must travel with each delegation event. Every HDP token encodes the scope, duration, and constraints of the human's delegation decision, and cryptographic verification is performed at each chain traversal \cite{hdp2025ietf}. However, HDP does not structurally enforce the integrity of the chain itself---it assumes the chain through which its tokens travel is trustworthy. In a mutable-chain environment, a compromised or misbehaving agent can present a false parent identity, inserting a false ancestor into the chain before presenting it to HDP token verification. The HDP token will verify correctly against the false chain because the token's own cryptographic integrity is independent of whether the chain topology is authentic \cite{hdp2025arxiv}. The combination of HDP-style token provenance with a structurally immutable parent chain closes both gaps simultaneously: HDP provides per-decision human authorization, and RSP provides the structurally immutable chain substrate that HDP's verification assumes but does not itself enforce.

% ---------------------------------------------------------------
\subsection{Accountability and Provenance in Distributed AI Architectures}\label{sec:provenance}

The challenge of tracing actions to their originating authority in distributed AI systems has been treated in the literature primarily as a provenance problem. Provenance refers to the complete record of the causal chain connecting an output to its input sources and the agents that transformed those sources along the way \cite{on_chain_provenance_ai}. The goal is to answer: \emph{who is responsible for this output, and through what chain of delegations did it come to pass?} In conventional distributed systems, provenance is achieved through append-only logging: each event is hash-linked to its predecessor, making retroactive modification computationally infeasible without detection \cite{blockchain_ai_governance}.

ChainScore Labs argues that only cryptographically verifiable, on-chain provenance for training data, model weights, and inference outputs can create the immutable audit trail required for safety, compliance, and trust in autonomous AI systems \cite{on_chain_provenance_ai}. Running model inference with zero-knowledge proofs (ZKML) creates an immutable, time-stamped record for every AI decision, transforming AI audits from ``opaque promises into verifiable records of model lineage and data origin'' \cite{on_chain_provenance_ai}. The authors project a 50\% improvement over traditional audits in simulation, with applicability across finance, healthcare, legal, and autonomous systems \cite{on_chain_provenance_ai}. FlexBlok proposes that blockchain-based distributed ledgers can establish independently verifiable, tamper-evident records for AI governance---recording model lifecycle events, logging AI decisions, and managing digital identities---without requiring trust in any single operator \cite{blockchain_ai_governance}. The approach integrates cryptographic verification, zero-knowledge proofs, and federated logging to balance verifiability with data confidentiality \cite{blockchain_ai_governance}.

The Zylos accountability analysis identifies three structural properties that mature accountability requires: signed identity (logs without signed identity cannot be verified), delegation chains (identity without delegation chains is incomplete), and policy enforcement (delegation chains without policy enforcement are unenforceable) \cite{zylos2026}. The analysis finds that among agents surveyed by the 2025 MIT AI Agent Index, only one (ChatGPT Agent) was found to use cryptographic request signing---suggesting that even prominent deployments largely lack standardized audit logging, identity verification, or delegation chain tracing \cite{zylos2026}. The analysis further establishes that accountability in multi-agent AI is not primarily a logging problem but an identity and authority problem \cite{zylos2026}. This reframing is significant: it shifts the design question from ``what should we log?'' to ``what structural invariants must hold for accountability to be preserved?''

The Warrant Certificate Authorities (WCA) framework addresses a complementary problem: certifying data provenance as data moves across tool-call boundaries in LLM-based agent systems \cite{wca2025}. WCA builds an end-to-end provenance layer with complete mediation, tamperproofness, and verifiability, drawing on PKI, operating system provenance mechanisms (IMA, CamFlow, LPM), and supply-chain security frameworks (SLSA, in-toto) \cite{wca2025}. The framework introduces Warrant Attestation Levels (WAL-0 to WAL-3) as a staged adoption model similar to SLSA build levels, providing a pathway from ad-hoc provenance to rigorously verified data lineage \cite{wca2025}. WAL-3 represents the highest assurance level, requiring cryptographic verification at every tool-call boundary with formal proof of provenance continuity \cite{wca2025}.

MAIF (Multimodal Artifact File Format) introduces an artifact-centric paradigm that grounds agent behavior in persistent, verifiable data artifacts embedded with semantic representations, cryptographic provenance, and fine-grained access controls \cite{maif2025arxiv}. Rather than relying on ephemeral task records, MAIF turns data into an active enforcement mechanism, providing immutable audit trails and chain-of-custody for AI operations that satisfy regulatory requirements such as the EU AI Act \cite{maif2025arxiv}. Production-ready features include high-speed streaming, cross-modal attention, semantic compression with fidelity preservation, real-time tamper detection, and behavioral anomaly analysis \cite{maif2025arxiv}.

The common property across all these systems is that provenance must be structurally immutable, not merely recorded. Blockchain-based approaches achieve this through distributed consensus; WCA achieves it through PKI-based warrant certificates; MAIF achieves it through artifact-centric enforcement. Each approach trades operational complexity for different forms of tamper-evidence. The BX3 Framework's forensic ledger achieves the same provenance properties within its own execution environment: append-only storage, hash-chained entries, and independent verifiability through hash-chain recomputation without reliance on the system's own operators \cite{bx3framework2026}. This matches what blockchain-based AI provenance systems seek to provide, but with lower operational overhead and no dependency on external distributed consensus.

% ---------------------------------------------------------------
\subsection{Fixed vs. Mutable Delegation Chains}\label{sec:fixed-v-mutable}

The architectural choice between fixed and mutable delegation chains has consequences that ramify across every accountability property of a multi-agent system. A mutable delegation chain allows parent references to be updated during execution, providing genuine operational advantages: agents can be dynamically reassigned in response to workload shifts, chains can be rebalanced to respond to performance bottlenecks, and failed agents can be replaced by transparently redirecting their children to new parents. These conveniences are the reason that essentially all production multi-agent systems today implement mutable chains.

The accountability cost of this convenience is severe. When a parent reference can change after the fact, the chain of responsibility becomes history-dependent in a way that defeats audit. An action taken at time $t$ may be attributable to one parent, but if the reference is later updated to point to a different agent, the same action becomes attributable to a different parent retroactively. Reconstructing the authoritative attribution chain requires capturing the complete mutable state at every point in time---a requirement that most systems fail to meet. In practice, mutable-chain systems either do not capture this history or capture it at insufficient fidelity to support definitive attribution. The SentinelAgent framework identifies this as the core vulnerability: in mutable chains, accountability is transitive only if provenance is immutable, but the framework acknowledges that ``the chain topology itself is mutable and therefore untrusted'' \cite{sentinelagent2025arxiv}.

A fixed delegation chain establishes parent references at node creation and preserves them unchanged for the node's lifetime. This rigidity forgoes the operational flexibilities described above, but it produces a chain that is self-certifying: any participant can reconstruct the authoritative chain at any time by walking parent pointers from any node to the root. No runtime state capture is required, because the chain structure itself is the record. This property---that the chain of custody is inherent in the structure rather than stored separately---is what makes fixed chains suitable for high-stakes accountability environments where the record must be trustworthy independent of the system's own operators. The practical cost of fixed chains is real but manageable: load rebalancing requires spawning new nodes rather than redirecting existing ones; a failed agent cannot be transparently replaced by reassigning its children. When an agent fails in an RSP chain, its children remain attached to their original parent, and the failure is surfaced through the escalation path rather than concealed by a transparent reassignment that would break the provenance chain. In high-stakes environments---legal reasoning, medical diagnosis, financial compliance, critical infrastructure orchestration---temporary throughput reduction during a failure incident is far less costly than accountability loss.

Policy-based approaches to chain integrity (access-control lists, role-based security models, audit logs) share a fundamental vulnerability: they trust the agent to enforce constraints against its own operational interests \cite{nist2023}. An agent near completion of a valuable task has a strong incentive to suppress escalation, exceed its token budget, or silently absorb a failure rather than surface it. The ReDel toolkit provides event-driven logging and replay for recursive multi-agent systems, enabling post-hoc analysis of delegation graphs \cite{redel2025arxiv}. However, ReDel assumes the delegation topology is correct and does not structurally enforce it; its value is diagnostic rather than preventive. The SentinelAgent framework introduces a delegation chain calculus for runtime enforcement \cite{sentinelagent2025arxiv}, but its enforcement depends on the runtime correctly interpreting chain topology---a property that mutable chains cannot guarantee. RSP addresses this gap at its root: the parent pointer is made architecturally immutable by the BX3 Fact Layer, not a policy parameter that the agent can rewrite or a logging convention that the agent can bypass.

The contrast between mutable and fixed delegation chains is not merely operational but epistemological. In a mutable chain, the question ``who authorized this action?'' has a time-dependent answer that may change as the chain evolves. In a fixed chain, the same question has a single, permanent answer traceable to the root. This permanence is a prerequisite for any accountability claim that must survive adversarial review---legal discovery, regulatory audit, or post-incident investigation. The EU AI Act's high-risk AI system requirements, fully effective August 2026, mandate automatic logging of agent actions throughout the system lifecycle with logs retained for 10 years for critical applications \cite{eu_ai_act}. Fixed delegation chains provide the structural foundation that makes such logging meaningful: without an immutable chain topology, a decade-old log entry may no longer correspond to any trustworthy attribution chain.

% ---------------------------------------------------------------
\subsection{Hierarchical Task Decomposition in AI Systems}\label{sec:task-decomposition}

Hierarchical task decomposition is the technique by which a complex goal is recursively subdivided into sub-goals assignable to specialized agents, enabling multi-agent systems to achieve coverage of tasks that no single agent could complete. The technique has deep roots in classical AI planning: HTN (Hierarchical Task Network) planning formalizes task decomposition as a process of recursively refining high-level goals into primitive actions executable by planning operators \cite{htn1993,htn2000}. SHOP (Simple Hierarchical Ordered Planner) demonstrated that HTN planning could scale to real-world domains by decomposing tasks into hierarchical networks processed in a top-down, left-to-right order \cite{htn2000}. The key insight from classical HTN planning that remains relevant for modern LLM-based agents is that decomposition produces a structured search space: the quality of the final plan is constrained by the quality of the decomposition decisions at each level of the hierarchy \cite{htn1993,htn2000}.

AgentOrchestra provides a systematic framework for hierarchical multi-agent task solving in which a manager agent decomposes goals and assigns subtasks to specialists, iterating until quality criteria are satisfied \cite{agentorchestra2025}. The framework introduces a Planning Agent as central orchestrator dedicated to high-level reasoning, task decomposition, and adaptive coordination, alongside specialized sub-agents (Deep Researcher Agent, Browser Use Agent, Deep Analyzer Agent, MCP Manager Agent) \cite{agentorchestra2025}. The planning agent sits at the top of the hierarchy, decomposing complex objectives and coordinating modular sub-agents responsible for domain-specific processing and multimodal reasoning \cite{agentorchestra2025}. On the GAIA benchmark testing dataset, AgentOrchestra achieves an average score of 83.39\%, ranking among the top general-purpose agents \cite{agentorchestra2025}. IBM's analysis observes that ``relying on a single AI agent for a complex, multi-step task means the agent will typically use an expensive, large frontier AI model for every step, even when not every step requires this level of sophistication; agent hierarchies create the flexibility to fit the right AI model to each step'' \cite{ibm2025hierarchical}. This cost-efficiency argument alone has driven widespread adoption of hierarchical decomposition in production agent systems.

ReAcTree introduces a hierarchical task-planning framework that decomposes complex objectives into dynamically constructed agent trees of subgoals, where each node is managed by an LLM agent capable of reasoning, acting, and expanding the tree \cite{reactree2025arxiv}. Control-flow nodes coordinate overall strategy; a dual-memory system (episodic memory for subgoal-specific examples, working memory for environment-specific observations) supports context preservation across the decomposition hierarchy \cite{reactree2025arxiv}. ReAcTree demonstrates that controlled recursion in hierarchical agent trees enables long-horizon task planning that would otherwise fail due to context-length limits, achieving 61\% goal success on WAH-NL compared to 31\% for flat ReAct baselines \cite{reactree2025arxiv}. On the ALFRED benchmark, ReAcTree similarly outperforms strong baselines, demonstrating that the hierarchical decomposition approach generalizes across embodied AI task domains \cite{reactree2025arxiv}.

HiPlan extends the hierarchical paradigm further, introducing a framework that decouples high-level planning from execution in deep search tasks \cite{hiplan2025arxiv}. The system breaks complex queries into subtasks, delegates each to domain-specific AI agents equipped with external tools, and coordinates results through a structured integration mechanism \cite{hiplan2025arxiv}. Prompt optimization using TextGrad-inspired textual gradient updates refines each agent's prompts when planning fails or yields infeasible actions, improving future plan quality iteratively \cite{hiplan2025arxiv}. Evaluated on the MAT-THOR benchmark, HiPlan achieves 0.95 success on compound tasks, 0.84 on complex tasks, and 0.60 on vague tasks, surpassing the prior state-of-the-art by 2, 7, and 15 percentage points respectively \cite{hiplan2025arxiv}. Ablation studies show the hierarchical structure contributes approximately +59 percentage points to overall success, with prompt optimization and meta-prompt sharing contributing +37 and +4 respectively \cite{hiplan2025arxiv}.

The Recursive Delegation Engine (RDE) formalizes hierarchical decomposition as a tree-growing process driven by utility-based decision policies \cite{rde2025arxiv}. At each step, an agent decides whether to execute directly, decompose further, or delegate to another agent, forming a delegation topology modeled as a directed graph $G = (A, E)$ where agents are nodes and edges represent permissible delegation paths \cite{rde2025arxiv}. RDE employs multi-armed bandit algorithms (epsilon-greedy, UCB/Beta-UCB, Thompson Sampling) to optimally select delegation paths based on observed success rates, with binary rewards propagating back through the chain to update per-agent performance counters \cite{rde2025arxiv}. This provides a principled mechanism for adaptive routing of tasks through agent hierarchies with formal guarantees on exploration-exploitation tradeoffs \cite{rde2025arxiv}.

Galileo AI's survey of multi-agent coordination failures identifies three emergent risks that appear specifically when delegation chains grow deep without constraint \cite{galileo2025survey}. First, goal coherence degrades as context diverges from original intent at each level of delegation. Second, attribution ambiguity makes it impossible to trace outputs to specific decomposition decisions. Third, system paralysis occurs when the coordination overhead of managing a deeply nested chain exceeds the system's operational capacity \cite{galileo2025survey}. These are not design flaws in any particular framework but emergent properties of recursive decomposition without structural depth constraints. The ReCAP framework introduces recursive context-aware reasoning and planning for long-horizon benchmarks \cite{recap2025arxiv}, acknowledging that managing context propagation across deep decomposition hierarchies is an unsolved problem.

Existing approaches to depth management are uniformly policy-based. Timeouts, token budgets, maximum chain length parameters, and human review triggers are all configured as operational parameters in the agent runtime and enforced by the runtime's execution logic \cite{galileo2025survey}. The fundamental vulnerability of policy-based depth management is structural: a policy constraint can be relaxed by the agent subject to it, because the agent controls the execution path that checks the constraint. An agent near completion of a valuable task has a strong incentive to suppress escalation, exceed its depth budget, or silently absorb a failure rather than surface it to the human anchor and risk termination. RSP addresses this by making the depth bound $D_{\max}$ a constraint enforced by the BX3 Bounds Engine layer, which operates in a separate execution context from the agents it constrains, so that no agent can modify the enforcement logic even indirectly.

% ---------------------------------------------------------------
\subsection{The BX3 Framework as Substrate for Immutable Parent Chains}\label{sec:bx3-substrate}

The BX3 Framework is an accountability architecture for multi-agent AI systems that organizes every node into three immutable functional layers: the Purpose layer, the Bounds Engine, and the Fact Layer \cite{bx3framework2026,bx3framework,bailoutprotocol2026}. The Purpose layer defines authorization policy, establishes the human anchor $h(n)$ for each node, and receives all escalated exceptions. The Bounds Engine evaluates operating range constraints including the depth bound $D_{\max}$, the spawn necessity criterion, and resource allocation limits. The Fact Layer maintains the append-only forensic ledger, enforces the immutable parent pointer via pre-commit hooks, and executes the structural invariants that make all other constraints enforceable \cite{bx3framework,bailoutprotocol2026}.

The framework's five pillars form an interdependent accountability system whose correctness properties require the simultaneous operation of all five \cite{bx3framework}. Pillar 1 (Loop Isolation) ensures that no two nodes share mutable execution state without mediated communication through the Fact Layer, preventing state corruption and ensuring that the provenance of every decision is traceable. Pillar 2 (Recursive Spawning with Immutable Parent Pointers) establishes the parent reference as a structurally immutable property established at node creation. Pillar 3 (Spatial Firewall) prevents cross-domain state contamination between nodes operating in distinct domains. Pillar 4 (Sandbox Gate) enforces the operating range constraints of the Bounds Engine as a structural gate rather than a policy parameter. Pillar 5 (Bailout Protocol) closes the accountability loop by compulsorily escalating any unresolved exception to the human anchor $h(n)$ through the immutable parent chain, bypassing all intermediate machine actors \cite{bx3framework,bailoutprotocol2026}.

The immutable parent pointer is Pillar 2 of the BX3 Framework. It is not a policy convention, not an operational guideline, and not a trust convention. It is a structural requirement enforced by the Fact Layer at node creation via a pre-commit hook that seals the parent reference before the node's execution context is initialized. After this point, the reference is cryptographically bound to the node's identity and cannot be modified by any operation, including operations initiated by the node itself, by administrative actors, or by any other framework component. This structural immutability is the propagation substrate for the Bailout Protocol: exception states have a structurally fixed escalation path through the immutable parent chain that no machine actor can redirect \cite{bailoutprotocol2026}.

The three-layer architecture is the minimal structure that provides RSP's correctness guarantees. Removing the Purpose layer eliminates the exclusive terminus of escalation, allowing machine actors to absorb exceptions and issue terminal directives, causing the chain to terminate in autonomous resolution rather than human judgment. Removing the Bounds Engine converts the depth bound from a structural enforcement into a policy parameter that operational urgency can exceed. Removing the Fact Layer makes the immutable parent pointer unenforceable; the parent reference becomes mutable and the forensic ledger cannot be independently trusted. This three-layer minimality is formally established in the BX3 Framework's integration specification for Recursive Spawning, which proves that no proper subset of layers achieves the same correctness guarantees \cite{bx3framework,bailoutprotocol2026}. The property is not empirical but structural: each guarantee depends on the causal mechanism that only that specific layer provides.

The BX3 forensic ledger provides the append-only, hash-chained record of all spawn events, escalation events, and resolution directives that is the prerequisite for independently verifiable provenance. Maintained by the Fact Layer and independently verifiable through hash-chain recomputation, the ledger requires no trusted operator and no public blockchain infrastructure. Its properties match what blockchain-based AI provenance systems seek to provide \cite{on_chain_provenance_ai,blockchain_ai_governance}, but achieved within the framework's own execution environment with lower operational overhead and no dependency on external distributed consensus. This makes the ledger suitable for regulatory audit, legal discovery, and adversarial review in high-stakes deployment environments \cite{bx3framework2026}.

The sovereign agent identity specification extends the BX3 Framework's accountability architecture to interoperable agent profiles, establishing canonical naming conventions, ontological scope boundaries, and cryptographic identity guarantees that enable trustworthy delegation across organizational boundaries \cite{sovereignagent}. This interoperability ensures that immutable parent chains can be verified not only within a single BX3 deployment but across different organizational contexts that share the same identity and provenance conventions \cite{sovereignagent}. The specification aligns with the IPSIE (Interoperable Profiles for Secure Identity and Entities) working group efforts to standardize agent identity for autonomous AI systems \cite{sovereignagent}.

% ---------------------------------------------------------------
% References
\bibliographystyle{plain}
\bibliography{RecursiveSpawning_Section2_Background_v7}

\end{document}
